% =====================================================================
\chapter{The base paper at a glance}\label{ch:p1paper}
% =====================================================================
\section{Motivation given by the paper}
The paper starts from three observations (Section~I of~\cite{tan2023}):
\begin{enumerate}[leftmargin=*]
\item Vehicle networks (IoVs) normally rely on roadside units (RSUs) at the edge. Existing authentication schemes assume these RSUs work. In a disaster
they may not, and then vehicle-to-infrastructure (V2I) communication fails.
\item Modern roads are three-dimensional (bridges, flyovers, tunnels), so flat ground coverage is not enough. Satellites have high latency and cost, so
the paper prefers UAVs as an aerial extension of the network.
\item Ordinary UAVs fly for less than an hour on battery. A \textbf{tethered UAV} gets power through a cable from a ground vehicle, so it can hover for a
long time and act as a mobile base station and edge server.
\end{enumerate}

\begin{figure}[H]
\centering
\resizebox{0.92\linewidth}{!}{%
\begin{tikzpicture}[font=\sffamily\small]
  \node[card=deep,text width=5.6cm,minimum height=4.0cm] (a) at (0,0) {};
  \node[font=\sffamily\bfseries,text=deep] at (0,1.65) {Normal IoV};
  \pic[color=deep,scale=1.1] at (0,0.35) {tower};
  \foreach \x in {-1.6,-0.5,0.6,1.7}{\pic[color=vehc,scale=0.8] at (\x,-0.9) {car};}
  \node[font=\scriptsize\sffamily,text=deep] at (0,-1.6) {vehicles talk through the RSU};
  \node[card=bad,text width=5.6cm,minimum height=4.0cm] (b) at (7,0) {};
  \node[font=\sffamily\bfseries,text=bad] at (7,1.65) {After a disaster};
  \pic[color=black!45,scale=1.1] at (7,0.35) {tower};
  \draw[bad,line width=2pt] (6.6,-0.15) -- (7.4,0.85) (6.6,0.85) -- (7.4,-0.15);
  \foreach \x in {5.4,6.5,7.6,8.7}{\pic[color=vehc,scale=0.8] at (\x,-0.9) {car};}
  \node[font=\scriptsize\sffamily,text=bad] at (7,-1.6) {RSU down: no V2I link};
  \node[card=hoc,text width=5.6cm,minimum height=4.0cm] (c) at (14,0) {};
  \node[font=\sffamily\bfseries,text=hoc] at (14,1.65) {The paper's idea};
  \pic[color=hoc,scale=0.9] at (13.3,0.75) {drone};
  \draw[tether,line width=1.4pt] (13.3,0.4) -- (13.3,-0.55);
  \pic[color=hoc,scale=0.9] at (13.3,-0.8) {truck};
  \pic[color=uavc,scale=0.6] at (14.6,0.8) {drone};
  \pic[color=uavc,scale=0.6] at (15.4,0.3) {drone};
  \node[font=\scriptsize\sffamily,text=hoc] at (14,-1.6) {TUAV on a tether + UAVs};
  \draw[-{Stealth[length=3mm]},very thick,black!50] (a.east) -- (b.west);
  \draw[-{Stealth[length=3mm]},very thick,black!50] (b.east) -- (c.west);
\end{tikzpicture}}
\caption{Motivation of the base paper. Left: in a normal IoV, vehicles communicate through a roadside unit (RSU). Middle: after a disaster the RSU is down and the vehicle-to-infrastructure link is lost. Right: the paper's idea -- a tethered UAV (TUAV), powered through a cable by a ground vehicle (ECRV), acts as the base station, and battery UAVs extend its coverage.}
\label{fig:motivation}
\end{figure}

\section{Contributions claimed by the paper}
\begin{claimbox}
\begin{itemize}[leftmargin=*]
\item ``Flexible and dynamic TUAV-enabled edge architecture for infrastructure-less IoV'' -- the TUAV replaces the failed edge server; the authors state
that this is the first use of a TUAV-enabled edge network for infrastructure-less IoV.
\item ``Secure and efficient TUAV-assisted group authentication mechanism with anonymity'' -- a certificateless group authentication with batch
validation, where ``the complex pairing computation tasks are allocated to the cloud server''.
\item ``Dynamic UAV management with updating strategy'' -- a group key built with the Chinese Remainder Theorem, with low cost for updating.
\end{itemize}
\hfill (Section I-A)
\end{claimbox}

\section{Related work discussed in the paper}
Section~II of the paper groups earlier work into a few families:
\begin{center}
\small
\begin{tabularx}{\linewidth}{@{}p{4.2cm} X@{}}
\toprule
\textbf{Family} & \textbf{Examples cited by the paper}\\
\midrule
Authenticated key management for IoV & dual authentication and key agreement for V2V; certificateless conditional privacy-preserving authentication
(CL-CPPA); fog-computing based AKM\\
Identity- and attribute-based schemes & pairing-free identity-based anonymous authentication with vehicle revocation; attribute-based encryption for
vehicular social networks\\
Lattice-based schemes & conditional privacy-preserving authentication based on the short integer solution problem\\
Group / batch authentication & batch authentication protocols that verify many vehicles at once, and certificateless aggregate signatures\\
UAV-assisted IoV & schemes that still need a working local edge infrastructure\\
\bottomrule
\end{tabularx}
\end{center}
The paper's main argument is that none of these works independently of the local edge infrastructure, which is the gap it addresses.

\section{Entities of the system}
\begin{center}
\small
\begin{tabularx}{\linewidth}{@{}l X@{}}
\toprule
\textbf{Entity} & \textbf{Role in the paper (Section III-D)}\\
\midrule
TA (trusted authority) & trusted server with large storage; does the setup and the user verification; reached through satellite (``space'') networks\\
TUAV (tethered UAV) & multirotor drone powered through a cable; acts as a substitute edge base station\\
ECRV & emergency communication response vehicle (a fire truck); powers the TUAV and acts as a control centre\\
UAV network & battery drones (under one hour of flight) that give 3D coverage\\
Vehicular network & vehicles that exchange driving data through the UAV network\\
\bottomrule
\end{tabularx}
\end{center}

\section{Security objectives stated by the paper}
The paper lists six goals (Section~III-E): \textbf{mutual authentication}, \textbf{user anonymity}, \textbf{unforgeability}, \textbf{non-repudiation},
\textbf{conditional privacy preservation} (only a law-enforcement agency can reveal a malicious user) and \textbf{session key establishment}.
Chapter~\ref{ch:base} shows how the protocol tries to reach them, and Chapter~\ref{ch:p1analysis} summarises how the paper argues it does.
